% --- Archon physics formalization source begin ---
% archon:physics
% archon:covers PhyXMiniProblems/problem_phyx_mini_0447.lean
% archon:source-report reports/phyx_mini/problem_phyx_mini_0447.source.json

\chapter{Physics problem phyx\_mini\_0447}
\label{ch:PhyXMiniProblems_problem_phyx_mini_0447}

\paragraph{Problem source.}
\#\# Physical scenario

A piston/cylinder arrangement, shown in figure, contains air at \$250 \textbackslash{}, \textbackslash{}text\{kPa\}\$ and \$300 \textbackslash{}, \textasciicircum{}\{\textbackslash{}circ\}\textbackslash{}text\{C\}\$. The \$50 \textbackslash{}, \textbackslash{}text\{kg\}\$ piston has a diameter of \$0.1 \textbackslash{}, \textbackslash{}text\{m\}\$ and initially pushes against the stops. The atmosphere is at \$100 \textbackslash{}, \textbackslash{}text\{kPa\}\$ and \$20 \textbackslash{}, \textasciicircum{}\{\textbackslash{}circ\}\textbackslash{}text\{C\}\$. The cylinder now cools as heat is transferred to the ambient surroundings.

\#\# Question

How far has the piston dropped when the temperature reaches ambient?

\#\# Answer choices

- A: A: 0.25 cm

- B: B: 7.8 cm

- C: C: 3.2 cm

- D: D: 5.3 cm

\#\# Figure caption (auxiliary; use the image as primary evidence)

The image depicts a cross-sectional diagram of a container filled with air. The container has a light blue region labeled "Air," representing the air inside. Above the air, there's a horizontal section that appears to be a piston, which is not labeled but shown with parallel lines indicating its presence and movement direction.

At the top of the image, the label \textbackslash{}( P\_0 \textbackslash{}) represents atmospheric pressure. An arrow labeled \textbackslash{}( g \textbackslash{}) points downward outside the container, indicating the direction of gravitational force. To the right, a vertical arrow specifies a height of 25 cm, likely indicating the height of the air column in the container.

The container is shown with thick gray walls. There are small gray blocks on either side of the top section, possibly representing support structures or seals.

\#\# Dataset metadata

Ideal Gases and Kinetic Theory; Multi-Formula Reasoning, Physical Model Grounding Reasoning

\paragraph{Recorded answer/context.}
D: D: 5.3 cm

\paragraph{Figure/image path.}
/root/proposal\_for\_physic/science-mango/phyx\_mini\_run/phyx\_data/test\_image/447.png

\paragraph{Formalization target.}
create a compiling Lean file with sorry bodies at `PhyXMiniProblems/problem\_phyx\_mini\_0447.lean`. The Lean declarations must preserve the physical quantities, dimensions or dimensional roles, figure labels, governing-law hypotheses, and final relation expressed by this problem.
Use Mathlib/Physlib names found through LeanExplore where available. If a physics API is missing, introduce faithful local abstractions rather than scalar placeholder aliases.

\begin{theorem}[Physics formalization target]
\label{thm:physics:phyx_mini_0447:target}
\lean{PhyXMiniProblems.ProblemPhyXMini0447.problem_phyx_mini_0447}
\uses{def:physics:phyx-mini-0447:phyxminiproblems-problemphyxmini0447-lengthincentimeters, def:physics:phyx-mini-0447:phyxminiproblems-problemphyxmini0447-coolingpistonsetup, def:physics:phyx-mini-0447:phyxminiproblems-problemphyxmini0447-matchescoolingpistonscenario, def:physics:phyx-mini-0447:phyxminiproblems-problemphyxmini0447-matchesproblemreadouts, def:physics:phyx-mini-0447:phyxminiproblems-problemphyxmini0447-matchessuppliedpistoncylinderfigure, def:physics:phyx-mini-0447:phyxminiproblems-problemphyxmini0447-usestextbookphysicalconstants, def:physics:phyx-mini-0447:phyxminiproblems-problemphyxmini0447-hasphysicalcoolingpistonparameters, def:physics:phyx-mini-0447:phyxminiproblems-problemphyxmini0447-reachesambienttemperature, def:physics:phyx-mini-0447:phyxminiproblems-problemphyxmini0447-satisfiescoolingpistonlaws, def:physics:phyx-mini-0447:phyxminiproblems-problemphyxmini0447-displayeddropcentimeters, def:physics:phyx-mini-0447:phyxminiproblems-problemphyxmini0447-recordeddatasetanswer, def:physics:phyx-mini-0447:phyxminiproblems-problemphyxmini0447-roundstonearesttenthcentimeter, def:physics:phyx-mini-0447:phyxminiproblems-problemphyxmini0447-isuniqueclosestchoice}
The assigned autoformalize agent should translate this physics problem into Lean declarations in the covered file, with theorem and lemma proof bodies written as `by sorry`.
\end{theorem}
\begin{proof}
This is an autoformalization task, not a proof task. Produce statements that can later be proved by the physics prover without weakening the physical model.
\end{proof}

% --- Archon declaration topology begin ---
\section{Declaration topology}
\begin{definition}[Length Quantity]
  \label{def:physics:phyx-mini-0447:phyxminiproblems-problemphyxmini0447-lengthquantity}
  \lean{PhyXMiniProblems.ProblemPhyXMini0447.LengthQuantity}
  A nonnegative physical length, independent of the chosen unit system
\end{definition}
\begin{proof}
  This is the declared model definition of Length Quantity.
\end{proof}
\begin{definition}[Volume Quantity]
  \label{def:physics:phyx-mini-0447:phyxminiproblems-problemphyxmini0447-volumequantity}
  \lean{PhyXMiniProblems.ProblemPhyXMini0447.VolumeQuantity}
  A nonnegative physical volume, with dimension L³
\end{definition}
\begin{proof}
  This is the declared model definition of Volume Quantity.
\end{proof}
\begin{definition}[Mass Quantity]
  \label{def:physics:phyx-mini-0447:phyxminiproblems-problemphyxmini0447-massquantity}
  \lean{PhyXMiniProblems.ProblemPhyXMini0447.MassQuantity}
  A nonnegative physical mass
\end{definition}
\begin{proof}
  This is the declared model definition of Mass Quantity.
\end{proof}
\begin{definition}[Acceleration Quantity]
  \label{def:physics:phyx-mini-0447:phyxminiproblems-problemphyxmini0447-accelerationquantity}
  \lean{PhyXMiniProblems.ProblemPhyXMini0447.AccelerationQuantity}
  A nonnegative acceleration magnitude, with dimension L T⁻²
\end{definition}
\begin{proof}
  This is the declared model definition of Acceleration Quantity.
\end{proof}
\begin{definition}[Force Quantity]
  \label{def:physics:phyx-mini-0447:phyxminiproblems-problemphyxmini0447-forcequantity}
  \lean{PhyXMiniProblems.ProblemPhyXMini0447.ForceQuantity}
  A nonnegative force magnitude, with dimension M L T⁻²
\end{definition}
\begin{proof}
  This is the declared model definition of Force Quantity.
\end{proof}
\begin{definition}[length In Meters]
  \label{def:physics:phyx-mini-0447:phyxminiproblems-problemphyxmini0447-lengthinmeters}
  \lean{PhyXMiniProblems.ProblemPhyXMini0447.lengthInMeters}
  \uses{def:physics:phyx-mini-0447:phyxminiproblems-problemphyxmini0447-lengthquantity}
  Read a physical length in metres
\end{definition}
\begin{proof}
  This is the declared model definition of length In Meters.
\end{proof}
\begin{definition}[length In Centimeters]
  \label{def:physics:phyx-mini-0447:phyxminiproblems-problemphyxmini0447-lengthincentimeters}
  \lean{PhyXMiniProblems.ProblemPhyXMini0447.lengthInCentimeters}
  \uses{def:physics:phyx-mini-0447:phyxminiproblems-problemphyxmini0447-lengthquantity, def:physics:phyx-mini-0447:phyxminiproblems-problemphyxmini0447-lengthinmeters}
  Read a physical length in centimetres
\end{definition}
\begin{proof}
  This is the declared model definition of length In Centimeters.
\end{proof}
\begin{definition}[area In Square Meters]
  \label{def:physics:phyx-mini-0447:phyxminiproblems-problemphyxmini0447-areainsquaremeters}
  \lean{PhyXMiniProblems.ProblemPhyXMini0447.areaInSquareMeters}
  Read a physical area in square metres
\end{definition}
\begin{proof}
  This is the declared model definition of area In Square Meters.
\end{proof}
\begin{definition}[volume In Cubic Meters]
  \label{def:physics:phyx-mini-0447:phyxminiproblems-problemphyxmini0447-volumeincubicmeters}
  \lean{PhyXMiniProblems.ProblemPhyXMini0447.volumeInCubicMeters}
  \uses{def:physics:phyx-mini-0447:phyxminiproblems-problemphyxmini0447-volumequantity}
  Read a physical volume in cubic metres
\end{definition}
\begin{proof}
  This is the declared model definition of volume In Cubic Meters.
\end{proof}
\begin{definition}[mass In Kilograms]
  \label{def:physics:phyx-mini-0447:phyxminiproblems-problemphyxmini0447-massinkilograms}
  \lean{PhyXMiniProblems.ProblemPhyXMini0447.massInKilograms}
  \uses{def:physics:phyx-mini-0447:phyxminiproblems-problemphyxmini0447-massquantity}
  Read a physical mass in kilograms
\end{definition}
\begin{proof}
  This is the declared model definition of mass In Kilograms.
\end{proof}
\begin{definition}[acceleration In Meters Per Second Squared]
  \label{def:physics:phyx-mini-0447:phyxminiproblems-problemphyxmini0447-accelerationinmeterspersecondsquared}
  \lean{PhyXMiniProblems.ProblemPhyXMini0447.accelerationInMetersPerSecondSquared}
  \uses{def:physics:phyx-mini-0447:phyxminiproblems-problemphyxmini0447-accelerationquantity}
  Read an acceleration magnitude in metres per second squared
\end{definition}
\begin{proof}
  This is the declared model definition of acceleration In Meters Per Second Squared.
\end{proof}
\begin{definition}[force In Newtons]
  \label{def:physics:phyx-mini-0447:phyxminiproblems-problemphyxmini0447-forceinnewtons}
  \lean{PhyXMiniProblems.ProblemPhyXMini0447.forceInNewtons}
  \uses{def:physics:phyx-mini-0447:phyxminiproblems-problemphyxmini0447-forcequantity}
  Read a force magnitude in newtons
\end{definition}
\begin{proof}
  This is the declared model definition of force In Newtons.
\end{proof}
\begin{definition}[pressure In Pascals]
  \label{def:physics:phyx-mini-0447:phyxminiproblems-problemphyxmini0447-pressureinpascals}
  \lean{PhyXMiniProblems.ProblemPhyXMini0447.pressureInPascals}
  Read a Physlib pressure in pascals
\end{definition}
\begin{proof}
  This is the declared model definition of pressure In Pascals.
\end{proof}
\begin{definition}[pressure In Kilopascals]
  \label{def:physics:phyx-mini-0447:phyxminiproblems-problemphyxmini0447-pressureinkilopascals}
  \lean{PhyXMiniProblems.ProblemPhyXMini0447.pressureInKilopascals}
  \uses{def:physics:phyx-mini-0447:phyxminiproblems-problemphyxmini0447-pressureinpascals}
  Read a Physlib pressure in kilopascals
\end{definition}
\begin{proof}
  This is the declared model definition of pressure In Kilopascals.
\end{proof}
\begin{definition}[temperature In Kelvin]
  \label{def:physics:phyx-mini-0447:phyxminiproblems-problemphyxmini0447-temperatureinkelvin}
  \lean{PhyXMiniProblems.ProblemPhyXMini0447.temperatureInKelvin}
  Read a Physlib absolute temperature in kelvin. The storage unit is explicit because Temperature stores a nonnegative value in an arbitrary zero-preserving temperature unit
\end{definition}
\begin{proof}
  This is the declared model definition of temperature In Kelvin.
\end{proof}
\begin{definition}[celsius Zero In Kelvin]
  \label{def:physics:phyx-mini-0447:phyxminiproblems-problemphyxmini0447-celsiuszeroinkelvin}
  \lean{PhyXMiniProblems.ProblemPhyXMini0447.celsiusZeroInKelvin}
  The exact Celsius-zero offset, 273.15 K
\end{definition}
\begin{proof}
  This is the declared model definition of celsius Zero In Kelvin.
\end{proof}
\begin{definition}[temperature In Degrees Celsius]
  \label{def:physics:phyx-mini-0447:phyxminiproblems-problemphyxmini0447-temperatureindegreescelsius}
  \lean{PhyXMiniProblems.ProblemPhyXMini0447.temperatureInDegreesCelsius}
  \uses{def:physics:phyx-mini-0447:phyxminiproblems-problemphyxmini0447-temperatureinkelvin, def:physics:phyx-mini-0447:phyxminiproblems-problemphyxmini0447-celsiuszeroinkelvin}
  Affine Celsius readout of a physical absolute temperature
\end{definition}
\begin{proof}
  This is the declared model definition of temperature In Degrees Celsius.
\end{proof}
\begin{definition}[Process State]
  \label{def:physics:phyx-mini-0447:phyxminiproblems-problemphyxmini0447-processstate}
  \lean{PhyXMiniProblems.ProblemPhyXMini0447.ProcessState}
  States separating the fixed-height and freely descending phases
\end{definition}
\begin{proof}
  This is the declared model definition of Process State.
\end{proof}
\begin{definition}[Enclosed Gas]
  \label{def:physics:phyx-mini-0447:phyxminiproblems-problemphyxmini0447-enclosedgas}
  \lean{PhyXMiniProblems.ProblemPhyXMini0447.EnclosedGas}
  The material confined below the piston
\end{definition}
\begin{proof}
  This is the declared model definition of Enclosed Gas.
\end{proof}
\begin{definition}[Cylinder Orientation]
  \label{def:physics:phyx-mini-0447:phyxminiproblems-problemphyxmini0447-cylinderorientation}
  \lean{PhyXMiniProblems.ProblemPhyXMini0447.CylinderOrientation}
  Orientation of the cylinder axis
\end{definition}
\begin{proof}
  This is the declared model definition of Cylinder Orientation.
\end{proof}
\begin{definition}[Cylinder Cross Section]
  \label{def:physics:phyx-mini-0447:phyxminiproblems-problemphyxmini0447-cylindercrosssection}
  \lean{PhyXMiniProblems.ProblemPhyXMini0447.CylinderCrossSection}
  Cross-sectional idealization used for the cylinder volume
\end{definition}
\begin{proof}
  This is the declared model definition of Cylinder Cross Section.
\end{proof}
\begin{definition}[Piston Support Regime]
  \label{def:physics:phyx-mini-0447:phyxminiproblems-problemphyxmini0447-pistonsupportregime}
  \lean{PhyXMiniProblems.ProblemPhyXMini0447.PistonSupportRegime}
  Mechanical support regime of the piston
\end{definition}
\begin{proof}
  This is the declared model definition of Piston Support Regime.
\end{proof}
\begin{definition}[Thermal Process]
  \label{def:physics:phyx-mini-0447:phyxminiproblems-problemphyxmini0447-thermalprocess}
  \lean{PhyXMiniProblems.ProblemPhyXMini0447.ThermalProcess}
  Thermal process described by the problem
\end{definition}
\begin{proof}
  This is the declared model definition of Thermal Process.
\end{proof}
\begin{definition}[Vertical Direction]
  \label{def:physics:phyx-mini-0447:phyxminiproblems-problemphyxmini0447-verticaldirection}
  \lean{PhyXMiniProblems.ProblemPhyXMini0447.VerticalDirection}
  Directions appearing in the vertical cross-sectional figure
\end{definition}
\begin{proof}
  This is the declared model definition of Vertical Direction.
\end{proof}
\begin{definition}[Figure Object]
  \label{def:physics:phyx-mini-0447:phyxminiproblems-problemphyxmini0447-figureobject}
  \lean{PhyXMiniProblems.ProblemPhyXMini0447.FigureObject}
  Physical objects visible in the supplied raster image
\end{definition}
\begin{proof}
  This is the declared model definition of Figure Object.
\end{proof}
\begin{definition}[Figure Label]
  \label{def:physics:phyx-mini-0447:phyxminiproblems-problemphyxmini0447-figurelabel}
  \lean{PhyXMiniProblems.ProblemPhyXMini0447.FigureLabel}
  Literal labels visible in the supplied raster image
\end{definition}
\begin{proof}
  This is the declared model definition of Figure Label.
\end{proof}
\begin{definition}[Piston Cylinder Figure]
  \label{def:physics:phyx-mini-0447:phyxminiproblems-problemphyxmini0447-pistoncylinderfigure}
  \lean{PhyXMiniProblems.ProblemPhyXMini0447.PistonCylinderFigure}
  \uses{def:physics:phyx-mini-0447:phyxminiproblems-problemphyxmini0447-verticaldirection, def:physics:phyx-mini-0447:phyxminiproblems-problemphyxmini0447-figureobject, def:physics:phyx-mini-0447:phyxminiproblems-problemphyxmini0447-figurelabel}
  Qualitative evidence transcribed from the primary image. Numerical figure evidence is connected to the initial thermodynamic state below, not encoded as a value of the requested displacement
\end{definition}
\begin{proof}
  This is the declared model definition of Piston Cylinder Figure.
\end{proof}
\begin{definition}[Gas State]
  \label{def:physics:phyx-mini-0447:phyxminiproblems-problemphyxmini0447-gasstate}
  \lean{PhyXMiniProblems.ProblemPhyXMini0447.GasState}
  \uses{def:physics:phyx-mini-0447:phyxminiproblems-problemphyxmini0447-lengthquantity, def:physics:phyx-mini-0447:phyxminiproblems-problemphyxmini0447-volumequantity}
  A gas state together with its cylinder-height observable
\end{definition}
\begin{proof}
  This is the declared model definition of Gas State.
\end{proof}
\begin{definition}[Cooling Piston Setup]
  \label{def:physics:phyx-mini-0447:phyxminiproblems-problemphyxmini0447-coolingpistonsetup}
  \lean{PhyXMiniProblems.ProblemPhyXMini0447.CoolingPistonSetup}
  \uses{def:physics:phyx-mini-0447:phyxminiproblems-problemphyxmini0447-lengthquantity, def:physics:phyx-mini-0447:phyxminiproblems-problemphyxmini0447-massquantity, def:physics:phyx-mini-0447:phyxminiproblems-problemphyxmini0447-accelerationquantity, def:physics:phyx-mini-0447:phyxminiproblems-problemphyxmini0447-forcequantity, def:physics:phyx-mini-0447:phyxminiproblems-problemphyxmini0447-processstate, def:physics:phyx-mini-0447:phyxminiproblems-problemphyxmini0447-enclosedgas, def:physics:phyx-mini-0447:phyxminiproblems-problemphyxmini0447-cylinderorientation, def:physics:phyx-mini-0447:phyxminiproblems-problemphyxmini0447-cylindercrosssection, def:physics:phyx-mini-0447:phyxminiproblems-problemphyxmini0447-pistonsupportregime, def:physics:phyx-mini-0447:phyxminiproblems-problemphyxmini0447-thermalprocess, def:physics:phyx-mini-0447:phyxminiproblems-problemphyxmini0447-pistoncylinderfigure, def:physics:phyx-mini-0447:phyxminiproblems-problemphyxmini0447-gasstate}
  Independent apparatus parameters and endpoint observables. The amount of air has an abstract physical carrier and only an explicitly calibrated mole readout. Neither pistonDrop nor the final height is assigned a numerical answer here
\end{definition}
\begin{proof}
  This is the declared model definition of Cooling Piston Setup.
\end{proof}
\begin{definition}[air Amount In Moles]
  \label{def:physics:phyx-mini-0447:phyxminiproblems-problemphyxmini0447-airamountinmoles}
  \lean{PhyXMiniProblems.ProblemPhyXMini0447.airAmountInMoles}
  \uses{def:physics:phyx-mini-0447:phyxminiproblems-problemphyxmini0447-coolingpistonsetup}
  Mole readout of the fixed amount of enclosed air
\end{definition}
\begin{proof}
  This is the declared model definition of air Amount In Moles.
\end{proof}
\begin{definition}[gas Temperature In Kelvin]
  \label{def:physics:phyx-mini-0447:phyxminiproblems-problemphyxmini0447-gastemperatureinkelvin}
  \lean{PhyXMiniProblems.ProblemPhyXMini0447.gasTemperatureInKelvin}
  \uses{def:physics:phyx-mini-0447:phyxminiproblems-problemphyxmini0447-temperatureinkelvin, def:physics:phyx-mini-0447:phyxminiproblems-problemphyxmini0447-processstate, def:physics:phyx-mini-0447:phyxminiproblems-problemphyxmini0447-coolingpistonsetup}
  Kelvin readout of the air temperature at a modeled state
\end{definition}
\begin{proof}
  This is the declared model definition of gas Temperature In Kelvin.
\end{proof}
\begin{definition}[gas Temperature In Degrees Celsius]
  \label{def:physics:phyx-mini-0447:phyxminiproblems-problemphyxmini0447-gastemperatureindegreescelsius}
  \lean{PhyXMiniProblems.ProblemPhyXMini0447.gasTemperatureInDegreesCelsius}
  \uses{def:physics:phyx-mini-0447:phyxminiproblems-problemphyxmini0447-temperatureindegreescelsius, def:physics:phyx-mini-0447:phyxminiproblems-problemphyxmini0447-processstate, def:physics:phyx-mini-0447:phyxminiproblems-problemphyxmini0447-coolingpistonsetup}
  Celsius readout of the air temperature at a modeled state
\end{definition}
\begin{proof}
  This is the declared model definition of gas Temperature In Degrees Celsius.
\end{proof}
\begin{definition}[ambient Temperature In Kelvin]
  \label{def:physics:phyx-mini-0447:phyxminiproblems-problemphyxmini0447-ambienttemperatureinkelvin}
  \lean{PhyXMiniProblems.ProblemPhyXMini0447.ambientTemperatureInKelvin}
  \uses{def:physics:phyx-mini-0447:phyxminiproblems-problemphyxmini0447-temperatureinkelvin, def:physics:phyx-mini-0447:phyxminiproblems-problemphyxmini0447-coolingpistonsetup}
  Kelvin readout of the ambient temperature
\end{definition}
\begin{proof}
  This is the declared model definition of ambient Temperature In Kelvin.
\end{proof}
\begin{definition}[ambient Temperature In Degrees Celsius]
  \label{def:physics:phyx-mini-0447:phyxminiproblems-problemphyxmini0447-ambienttemperatureindegreescelsius}
  \lean{PhyXMiniProblems.ProblemPhyXMini0447.ambientTemperatureInDegreesCelsius}
  \uses{def:physics:phyx-mini-0447:phyxminiproblems-problemphyxmini0447-temperatureindegreescelsius, def:physics:phyx-mini-0447:phyxminiproblems-problemphyxmini0447-coolingpistonsetup}
  Celsius readout of the ambient temperature
\end{definition}
\begin{proof}
  This is the declared model definition of ambient Temperature In Degrees Celsius.
\end{proof}
\begin{definition}[Matches Cooling Piston Scenario]
  \label{def:physics:phyx-mini-0447:phyxminiproblems-problemphyxmini0447-matchescoolingpistonscenario}
  \lean{PhyXMiniProblems.ProblemPhyXMini0447.MatchesCoolingPistonScenario}
  \uses{def:physics:phyx-mini-0447:phyxminiproblems-problemphyxmini0447-coolingpistonsetup}
  Qualitative physical conditions stated or implied by the problem
\end{definition}
\begin{proof}
  This is the declared model definition of Matches Cooling Piston Scenario.
\end{proof}
\begin{definition}[Matches Problem Readouts]
  \label{def:physics:phyx-mini-0447:phyxminiproblems-problemphyxmini0447-matchesproblemreadouts}
  \lean{PhyXMiniProblems.ProblemPhyXMini0447.MatchesProblemReadouts}
  \uses{def:physics:phyx-mini-0447:phyxminiproblems-problemphyxmini0447-lengthinmeters, def:physics:phyx-mini-0447:phyxminiproblems-problemphyxmini0447-massinkilograms, def:physics:phyx-mini-0447:phyxminiproblems-problemphyxmini0447-pressureinkilopascals, def:physics:phyx-mini-0447:phyxminiproblems-problemphyxmini0447-coolingpistonsetup, def:physics:phyx-mini-0447:phyxminiproblems-problemphyxmini0447-gastemperatureindegreescelsius, def:physics:phyx-mini-0447:phyxminiproblems-problemphyxmini0447-ambienttemperatureindegreescelsius}
  Numerical data stated in the prose. The requested drop and final gas-column height are deliberately absent
\end{definition}
\begin{proof}
  This is the declared model definition of Matches Problem Readouts.
\end{proof}
\begin{definition}[Matches Supplied Piston Cylinder Figure]
  \label{def:physics:phyx-mini-0447:phyxminiproblems-problemphyxmini0447-matchessuppliedpistoncylinderfigure}
  \lean{PhyXMiniProblems.ProblemPhyXMini0447.MatchesSuppliedPistonCylinderFigure}
  \uses{def:physics:phyx-mini-0447:phyxminiproblems-problemphyxmini0447-lengthincentimeters, def:physics:phyx-mini-0447:phyxminiproblems-problemphyxmini0447-figureobject, def:physics:phyx-mini-0447:phyxminiproblems-problemphyxmini0447-figurelabel, def:physics:phyx-mini-0447:phyxminiproblems-problemphyxmini0447-coolingpistonsetup}
  Primary-raster evidence: P₀ is above the piston, g points downward, air is below the piston, the piston touches the upper stops, and the height arrow reads 25 cm. The image contains no final displacement readout
\end{definition}
\begin{proof}
  This is the declared model definition of Matches Supplied Piston Cylinder Figure.
\end{proof}
\begin{definition}[Uses Textbook Physical Constants]
  \label{def:physics:phyx-mini-0447:phyxminiproblems-problemphyxmini0447-usestextbookphysicalconstants}
  \lean{PhyXMiniProblems.ProblemPhyXMini0447.UsesTextbookPhysicalConstants}
  \uses{def:physics:phyx-mini-0447:phyxminiproblems-problemphyxmini0447-accelerationinmeterspersecondsquared, def:physics:phyx-mini-0447:phyxminiproblems-problemphyxmini0447-coolingpistonsetup}
  Independent textbook constants used by the physical prediction
\end{definition}
\begin{proof}
  This is the declared model definition of Uses Textbook Physical Constants.
\end{proof}
\begin{definition}[Has Physical Cooling Piston Parameters]
  \label{def:physics:phyx-mini-0447:phyxminiproblems-problemphyxmini0447-hasphysicalcoolingpistonparameters}
  \lean{PhyXMiniProblems.ProblemPhyXMini0447.HasPhysicalCoolingPistonParameters}
  \uses{def:physics:phyx-mini-0447:phyxminiproblems-problemphyxmini0447-lengthinmeters, def:physics:phyx-mini-0447:phyxminiproblems-problemphyxmini0447-areainsquaremeters, def:physics:phyx-mini-0447:phyxminiproblems-problemphyxmini0447-volumeincubicmeters, def:physics:phyx-mini-0447:phyxminiproblems-problemphyxmini0447-massinkilograms, def:physics:phyx-mini-0447:phyxminiproblems-problemphyxmini0447-accelerationinmeterspersecondsquared, def:physics:phyx-mini-0447:phyxminiproblems-problemphyxmini0447-forceinnewtons, def:physics:phyx-mini-0447:phyxminiproblems-problemphyxmini0447-pressureinpascals, def:physics:phyx-mini-0447:phyxminiproblems-problemphyxmini0447-coolingpistonsetup, def:physics:phyx-mini-0447:phyxminiproblems-problemphyxmini0447-airamountinmoles, def:physics:phyx-mini-0447:phyxminiproblems-problemphyxmini0447-gastemperatureinkelvin, def:physics:phyx-mini-0447:phyxminiproblems-problemphyxmini0447-ambienttemperatureinkelvin}
  Positivity and nondegeneracy conditions selecting physical states
\end{definition}
\begin{proof}
  This is the declared model definition of Has Physical Cooling Piston Parameters.
\end{proof}
\begin{definition}[Reaches Ambient Temperature]
  \label{def:physics:phyx-mini-0447:phyxminiproblems-problemphyxmini0447-reachesambienttemperature}
  \lean{PhyXMiniProblems.ProblemPhyXMini0447.ReachesAmbientTemperature}
  \uses{def:physics:phyx-mini-0447:phyxminiproblems-problemphyxmini0447-coolingpistonsetup, def:physics:phyx-mini-0447:phyxminiproblems-problemphyxmini0447-gastemperatureinkelvin, def:physics:phyx-mini-0447:phyxminiproblems-problemphyxmini0447-ambienttemperatureinkelvin}
  The event named in the question: cooling continues until the air reaches the ambient absolute temperature. This fixes when displacement is measured, not the value of that displacement
\end{definition}
\begin{proof}
  This is the declared model definition of Reaches Ambient Temperature.
\end{proof}
\begin{definition}[Satisfies Cooling Piston Laws]
  \label{def:physics:phyx-mini-0447:phyxminiproblems-problemphyxmini0447-satisfiescoolingpistonlaws}
  \lean{PhyXMiniProblems.ProblemPhyXMini0447.SatisfiesCoolingPistonLaws}
  \uses{def:physics:phyx-mini-0447:phyxminiproblems-problemphyxmini0447-lengthinmeters, def:physics:phyx-mini-0447:phyxminiproblems-problemphyxmini0447-areainsquaremeters, def:physics:phyx-mini-0447:phyxminiproblems-problemphyxmini0447-volumeincubicmeters, def:physics:phyx-mini-0447:phyxminiproblems-problemphyxmini0447-massinkilograms, def:physics:phyx-mini-0447:phyxminiproblems-problemphyxmini0447-accelerationinmeterspersecondsquared, def:physics:phyx-mini-0447:phyxminiproblems-problemphyxmini0447-forceinnewtons, def:physics:phyx-mini-0447:phyxminiproblems-problemphyxmini0447-pressureinpascals, def:physics:phyx-mini-0447:phyxminiproblems-problemphyxmini0447-coolingpistonsetup, def:physics:phyx-mini-0447:phyxminiproblems-problemphyxmini0447-airamountinmoles, def:physics:phyx-mini-0447:phyxminiproblems-problemphyxmini0447-gastemperatureinkelvin}
  Governing geometry, mechanics, thermodynamics, and phase behavior: the piston face is circular and gas volume is area times column height; piston weight is m g; while the upper stops engage, the height remains fixed and the stops supply the excess downward reaction; after release, static force balance fixes gas pressure from atmospheric loading plus piston weight; the same closed amount of ideal air obeys P V = n R T at every state;
\end{definition}
\begin{proof}
  This is the declared model definition of Satisfies Cooling Piston Laws.
\end{proof}
\begin{theorem}[piston Drop Formula In Meters]
  \label{thm:physics:phyx-mini-0447:phyxminiproblems-problemphyxmini0447-pistondropformulainmeters}
  \lean{PhyXMiniProblems.ProblemPhyXMini0447.pistonDropFormulaInMeters}
  \uses{def:physics:phyx-mini-0447:phyxminiproblems-problemphyxmini0447-lengthinmeters, def:physics:phyx-mini-0447:phyxminiproblems-problemphyxmini0447-areainsquaremeters, def:physics:phyx-mini-0447:phyxminiproblems-problemphyxmini0447-forceinnewtons, def:physics:phyx-mini-0447:phyxminiproblems-problemphyxmini0447-pressureinpascals, def:physics:phyx-mini-0447:phyxminiproblems-problemphyxmini0447-coolingpistonsetup, def:physics:phyx-mini-0447:phyxminiproblems-problemphyxmini0447-gastemperatureinkelvin, def:physics:phyx-mini-0447:phyxminiproblems-problemphyxmini0447-ambienttemperatureinkelvin, def:physics:phyx-mini-0447:phyxminiproblems-problemphyxmini0447-hasphysicalcoolingpistonparameters, def:physics:phyx-mini-0447:phyxminiproblems-problemphyxmini0447-reachesambienttemperature, def:physics:phyx-mini-0447:phyxminiproblems-problemphyxmini0447-satisfiescoolingpistonlaws}
  Combining the endpoint ideal-gas equations, cylindrical geometry, final free-piston force balance, and the ambient-temperature event gives this general displacement relation. It contains no problem-specific answer value
\end{theorem}
\begin{proof}
  The conclusion follows from the governing relations and figure readouts named in the statement of piston Drop Formula In Meters.
\end{proof}
\begin{definition}[Answer Choice]
  \label{def:physics:phyx-mini-0447:phyxminiproblems-problemphyxmini0447-answerchoice}
  \lean{PhyXMiniProblems.ProblemPhyXMini0447.AnswerChoice}
  Labels of the four answers printed in the problem statement
\end{definition}
\begin{proof}
  This is the declared model definition of Answer Choice.
\end{proof}
\begin{definition}[displayed Drop Centimeters]
  \label{def:physics:phyx-mini-0447:phyxminiproblems-problemphyxmini0447-displayeddropcentimeters}
  \lean{PhyXMiniProblems.ProblemPhyXMini0447.displayedDropCentimeters}
  \uses{def:physics:phyx-mini-0447:phyxminiproblems-problemphyxmini0447-answerchoice}
  Piston drop in centimetres printed beside each answer label
\end{definition}
\begin{proof}
  This is the declared model definition of displayed Drop Centimeters.
\end{proof}
\begin{definition}[recorded Dataset Answer]
  \label{def:physics:phyx-mini-0447:phyxminiproblems-problemphyxmini0447-recordeddatasetanswer}
  \lean{PhyXMiniProblems.ProblemPhyXMini0447.recordedDatasetAnswer}
  \uses{def:physics:phyx-mini-0447:phyxminiproblems-problemphyxmini0447-answerchoice}
  Answer label recorded in the dataset metadata
\end{definition}
\begin{proof}
  This is the declared model definition of recorded Dataset Answer.
\end{proof}
\begin{definition}[Rounds To Nearest Tenth Centimeter]
  \label{def:physics:phyx-mini-0447:phyxminiproblems-problemphyxmini0447-roundstonearesttenthcentimeter}
  \lean{PhyXMiniProblems.ProblemPhyXMini0447.RoundsToNearestTenthCentimeter}
  The modeled drop rounds to a display given to the nearest 0.1 cm
\end{definition}
\begin{proof}
  This is the declared model definition of Rounds To Nearest Tenth Centimeter.
\end{proof}
\begin{definition}[Is Unique Closest Choice]
  \label{def:physics:phyx-mini-0447:phyxminiproblems-problemphyxmini0447-isuniqueclosestchoice}
  \lean{PhyXMiniProblems.ProblemPhyXMini0447.IsUniqueClosestChoice}
  \uses{def:physics:phyx-mini-0447:phyxminiproblems-problemphyxmini0447-answerchoice, def:physics:phyx-mini-0447:phyxminiproblems-problemphyxmini0447-displayeddropcentimeters}
  A chosen display is strictly closer than every alternative display
\end{definition}
\begin{proof}
  This is the declared model definition of Is Unique Closest Choice.
\end{proof}
% --- Archon declaration topology end ---
% --- Archon physics formalization source end ---
